\mode*
\section{Remotes and collaboration}
\label{sec:remotes}

The fourth educational objective is distribution: clones as full copies,
remotes as named other copies, push/pull as synchronisation---and the
collaboration workflows built on top (GitHub, pull requests).

\subsection{Documented misconceptions and difficulties}

Course observations document confusion about the distributed nature of the
system and about collaborating through a shared remote
\autocite{IsomottonenCochez2014}; a concrete instance is the
remote-tracking-name confusion:
\textcquote[Sec.~5.4]{IsomottonenCochez2014}{this `with and without prefix'
conventions confuses learners until they have used the system long enough and
grasped the need for specifying the remote context with a prefix}
(\texttt{origin/master} when listing, \texttt{origin master} when pushing).
The GitKit line catalogues misconceptions in the basic GitHub workflows
across institutions \autocite{Postner2026GitMisconceptions,Braught2024GitKit},
and the Shell Tutor covers Git in its shell-based tutoring
\autocite{Winder2024ShellTutor}.
Because the distributed state is what learners cannot see, visualising it is a
natural intervention: a recent split-view learning tool mirrors a partner's
actions on a shared remote in real time to build
\textcquote{Bucher2026GitTakesTwo}{awareness of distributed states,
coordination, and collaborative troubleshooting}.

\subsection{Candidate critical aspects}
\label{sec:remotes-aspects}

The distributed state is what learners cannot see, and two of the three
candidates rest on difficulties the literature documents; the first is
posited.

\begin{description}
  \item[Every clone is complete]
    There is no \enquote{main} copy by mechanism, only by convention;
    each clone holds the whole history.
    Posited from the object of learning; the GitKit catalogue's
    main-copy workflow misconceptions \autocite{Postner2026GitMisconceptions}
    point the same way, at poster weight.
    Tested by the item \emph{offline with a clone}
    (\cref{app:knowledge-items}).
  \item[Push and pull move commits]
    Synchronisation is between histories, not files; an uncommitted
    edit stays behind.
    Read off the remote-prefix confusion
    \autocite[Sec.~5.4]{IsomottonenCochez2014}---\texttt{origin/main} names
    a copy of a history, which a learner who moves files has no use
    for---and off the fact that where Git is a stated learning objective
    this is exactly what instructors test, with a pull that modifies
    the local directory, the remote, both or neither
    \autocite{Beckman2021GitLearningObjective}.
    Tested by the item \emph{what \texttt{git push} sends}.
  \item[Local commits are invisible to others until pushed]
    And theirs to you until pulled.
    Read off the documented confusion about the distributed state
    \autocite{IsomottonenCochez2014}, the hidden dependency between a
    local branch, its remote-tracking copy and the remote---a status that
    is \enquote{up-to-date with origin/master} while hundreds of commits
    behind \autocite{Church2014HiddenDependencies}---and the awareness a
    split-view tool sets out to build \autocite{Bucher2026GitTakesTwo}.
    Tested by the item \emph{why they cannot see it} and its open twin
    \emph{the teammate's clone} (\cref{app:open-items}).
\end{description}

\subsection{Patterns of variation}

\begin{itemize}
  \item Contrast (copy varies, history invariant): \texttt{git log} in
    the clone on the laptop and in the repository on GitHub show the
    same history; unplug the network, and the laptop's \texttt{git log}
    still does.
  \item Contrast (local vs remote): commit without pushing; observe the
    repository on GitHub unchanged; push; observe.
  \item Generalisation (workflow invariant, collaborators vary): the same
    add--commit--push cycle alone, then with a partner pulling.
  \item Fusion: a deliberate concurrent edit leading to a conflict,
    resolved while reading the graph.
\end{itemize}
